\Subhead{2.1.1}{Плоские однородные и неоднородные волны. Поток энергии. Поляризация}

Плоской называют волну, поверхности постоянной фазы которой — плоскости. Для монохроматического поля
\begin{equation}
 u(\bm r,t)=\operatorname{Re}\{U\exp[\ii\bm k\cdot\bm r-\ii\omega t]\},\qquad \bm k=\bm k'+\ii\bm k''.
\end{equation}
Фаза постоянна на $\bm k'\cdot\bm r=\mathrm{const}$, амплитуда — на $\bm k''\cdot\bm r=\mathrm{const}$. Если эти плоскости параллельны, волну называют однородной; неоднородная волна имеет непараллельные направления распространения фазы и затухания. Поглощающая однородная волна может иметь комплексный $k$: «однородная» не означает «незатухающая». Для эванесцентного поля вдоль границы фаза идёт вдоль поверхности, затухание — поперёк.

Векторное поле за период описывает в фиксированной точке отрезок, эллипс или окружность: линейная, эллиптическая или круговая поляризация. Для двух поперечных компонент равные амплитуды и фазовый сдвиг $\pm\pi/2$ дают круг. Поляризация описывает движение вектора, а не форму фазовой поверхности.

Поток энергии определяется природой поля: в акустике $\bm I=\langle p'\bm v\rangle$, в электродинамике $\bm S=\bm E\times\bm H$. У монохроматического поля средний поток содержит $\frac12\operatorname{Re}$ произведения одной амплитуды на сопряжённую другой. В анизотропной среде поток энергии может не совпадать с $\bm k'$; у эванесцентной волны мнимый поток характеризует реактивный обмен и не равен переносу средней мощности.
